% Source: https://artofproblemsolving.com/community/c3h2685387_does_a_linear_function_have_an_axis_of_symmetry
% Stable identity: aops-topic-2685387-q01
\textbf{(a)} If $L=\ell$, the reflection fixes every point of $\ell$. If $L\perp\ell$, let $X$ be any point of $\ell$. Its image $X'$ lies on the line through $X$ perpendicular to $L$, and that line is $\ell$ itself; so $X'\in\ell$, and the reflection maps $\ell$ onto $\ell$.

Conversely, suppose the reflection maps $\ell$ onto itself and $L\ne\ell$. Choose $X\in\ell$ not on $L$. Its image $X'$ is on $\ell$ and $X'\ne X$, so $\ell$ is the line $XX'$. But $XX'\perp L$ by the definition of reflection. Hence $\ell\perp L$. $\blacksquare$

\begin{asy}
usepackage("amsmath");
size(4.5cm);
defaultpen(fontsize(10pt));
pair P=(1,1);
draw((-2.5,-0.75)--(3.2,2.1),linewidth(0.9));
draw(P+1.2*dir(degrees((1,-2)))--P-1.1*dir(degrees((1,-2))),dashed);
pair X=(2.6,1.8),Xp=2*P-X;dot(P);dot(X);dot(Xp);
draw(X--Xp,dotted+gray);
pair v1=unit((2,1)),v2=unit((1,-2));draw(P+0.25*v1--P+0.25*(v1+v2)--P+0.25*v2);
label("$\ell$",(3.2,2.1),E);label("$L$",P+1.2*unit((1,-2)),S);
label("$P$",P,dir(72));label("$X$",X,SE);label("$X'$",Xp,NW);
\end{asy}

\textbf{(b)} The reflection in $x=h$ sends $(x,y)$ to $(2h-x,y)$. The graph is symmetric about $x=h$ exactly when $m(2h-x)+b=mx+b$ for all $x$, i.e. $2m(h-x)=0$ for all $x$, i.e. $m=0$. (This agrees with (a): a vertical line is perpendicular to $\ell$ only when $\ell$ is horizontal.) So $\boxed{m=0}$, and then \emph{every} vertical line is an axis of symmetry.
